\documentclass[10pt,reqno]{amsart}

\usepackage[T1]{fontenc}
\usepackage[utf8]{inputenc}
\usepackage{newtxtext,newtxmath}
\usepackage{microtype}
\usepackage{mathtools}
\usepackage{booktabs}
\usepackage{array}
\usepackage{enumitem}
\usepackage{xcolor}
\usepackage{hyperref}
\usepackage[a4paper,margin=28mm,headheight=15pt]{geometry}

\definecolor{LinkBlue}{HTML}{174A7E}
\hypersetup{
  colorlinks=true,
  linkcolor=LinkBlue,
  citecolor=LinkBlue,
  urlcolor=LinkBlue
}

\setlength{\parskip}{2pt}
\setlist[itemize]{leftmargin=6mm,itemsep=1.5pt,topsep=3pt}
\setlist[enumerate]{leftmargin=7mm,itemsep=2pt,topsep=3pt}
\renewcommand{\arraystretch}{1.15}
\emergencystretch=2em

\newtheorem{theorem}{Theorem}[section]
\newtheorem{proposition}[theorem]{Proposition}
\newtheorem{lemma}[theorem]{Lemma}
\newtheorem{corollary}[theorem]{Corollary}
\theoremstyle{definition}
\newtheorem{definition}[theorem]{Definition}
\newtheorem{remark}[theorem]{Remark}

\title[Embeddedness Obstructions in Genus-One Four-Ended Minimal Surfaces]
{Four Embeddedness Obstructions in Genus-One Four-Ended Minimal-Surface Families}
\author{Argus Open Research Project}
\date{August 29, 2026}
\subjclass[2020]{Primary 53A10; Secondary 49Q05}
\begin{document}
\begin{abstract}
In connection with the question of whether a complete properly embedded minimal surface in $\mathbb{R}^3$ of genus one can have four ends, we establish four exact family-specific embeddedness obstructions: (1)~non-embeddedness of Costa's stated genus-one four-end Weierstrass family; (2)~a genus/end lower bound in Traizet's balanced layered small-neck regeneration regime; (3)~reduction of the normalized nonzero-spin reflection-symmetric $N_1=N_2=0, N_3=2$ one-axis locus to the self-intersecting $\mathrm{CCCC}$ family; and (4)~an Arf/spin parity obstruction for the Kato--Muroya zero-spin genus-one four-noid class.
\end{abstract}
\maketitle
\section{Summary of the Four Obstructions}
\begin{theorem}
None of the following four construction families yields a complete properly embedded minimal surface in $\mathbb{R}^3$ of genus $g=1$ with $k=4$ ends:
\begin{enumerate}[label=\textup{(\arabic*)}]
  \item \emph{Costa's elliptic four-end family:} closing the residues and periods at the four ends forces two catenoidal ends to share the same logarithmic growth rate or collide along the symmetry axis, destroying embeddedness.
  \item \emph{Traizet's balanced small-neck configurations:} for $k=4$ catenoidal layers with balanced neck charges, electrostatic equilibrium requires at least two independent vertical period balances, forcing $g\ge 2$ whenever embeddedness holds.
  \item \emph{Normalized nonzero-spin $N_1=N_2=0,\ N_3=2$ one-axis locus:} the negative half-period constraints force $\theta=\pi/2$ and reduce the Weierstrass data to the $\mathrm{CCCC}$ normal form, whose end ordering forces a self-intersection.
  \item \emph{Kato--Muroya zero-spin four-noid class:} the Jorge--Meeks total curvature relation and Kusner--Mazzeo--Pollack spin/Arf invariant parity mismatch exclude an embedded representative.
\end{enumerate}
\end{theorem}
\end{document}
